% IGCSE 0478 blueprint demo -- proves the shared preamble's six custom
% environments compile and render distinctly. Not real content; never
% copy this file's body into 05_output, only its structure/pattern.
% IGCSE 0478 shared preamble -- structure only, NO topic content.
% \input this file as the FIRST line of every 05_output/*.tex (it includes
% \documentclass, so it must come before anything else). Defines the three
% custom environments every generated document must use exclusively.
% See CLAUDE.md "Strict Separation of Structure and Content".
\documentclass[11pt]{article}

\usepackage[a4paper, margin=2cm]{geometry}
\usepackage{xcolor}
\usepackage{tcolorbox}
\usepackage{enumitem}
\usepackage{amsmath}
\usepackage{hyperref}

\tcbuselibrary{skins, breakable}

% stratbox: an exam strategy / technique tip.
\newtcolorbox{stratbox}[1][Exam Strategy]{
  enhanced, breakable, colback=blue!5, colframe=blue!60!black,
  fonttitle=\bfseries, title=#1
}

% critbox: a common mistake / misconception warning.
\newtcolorbox{critbox}[1][Common Mistake]{
  enhanced, breakable, colback=red!5, colframe=red!60!black,
  fonttitle=\bfseries, title=#1
}

% scenario: an applied exam-style scenario or worked example.
\newtcolorbox{scenario}[1][Scenario]{
  enhanced, breakable, colback=green!5, colframe=green!40!black,
  fonttitle=\bfseries, title=#1
}


% --- Class-notes section skeleton -------------------------------------------
% One \subsection per syllabus subsection, e.g. 1.1, 1.2, 1.3 -- content
% comes from 03_knowledge_base/Topic_XX_Name.md.
%
% % IGCSE 0478 shared preamble -- structure only, NO topic content.
% \input this file as the FIRST line of every 05_output/*.tex (it includes
% \documentclass, so it must come before anything else). Defines the three
% custom environments every generated document must use exclusively.
% See CLAUDE.md "Strict Separation of Structure and Content".
\documentclass[11pt]{article}

\usepackage[a4paper, margin=2cm]{geometry}
\usepackage{xcolor}
\usepackage{tcolorbox}
\usepackage{enumitem}
\usepackage{amsmath}
\usepackage{hyperref}

\tcbuselibrary{skins, breakable}

% stratbox: an exam strategy / technique tip.
\newtcolorbox{stratbox}[1][Exam Strategy]{
  enhanced, breakable, colback=blue!5, colframe=blue!60!black,
  fonttitle=\bfseries, title=#1
}

% critbox: a common mistake / misconception warning.
\newtcolorbox{critbox}[1][Common Mistake]{
  enhanced, breakable, colback=red!5, colframe=red!60!black,
  fonttitle=\bfseries, title=#1
}

% scenario: an applied exam-style scenario or worked example.
\newtcolorbox{scenario}[1][Scenario]{
  enhanced, breakable, colback=green!5, colframe=green!40!black,
  fonttitle=\bfseries, title=#1
}

% \begin{document}
% \title{Topic XX: <Topic Name>}
% \maketitle
% \section*{<Topic Name>}
%   \subsection*{<Syllabus subsection, e.g. 1.1 Number systems>}
%     <content from 03_knowledge_base, as prose/itemize>
%     \begin{stratbox} ... \end{stratbox}   % optional, only if KB has a tip
%     \begin{critbox} ... \end{critbox}     % optional, only if KB flags a mistake
%     \begin{scenario} ... \end{scenario}   % optional, only if KB has a worked example
%   % TODO: Insufficient source content      <- used verbatim when KB has a gap
% \end{document}

% --- Worksheet section skeleton ---------------------------------------------
% Shares the Topic_XX_Name stem with its class notes (see Topic_XX_Name.tex
% vs. Topic_XX_Name_Worksheet.tex). Content comes from the linked QP/MS
% Markdown in 03_knowledge_base/PastPaper_<code>.md (## Question <label> /
% **Question:** ... / **Mark scheme:** ...) -- the question text and the
% mark scheme are two separate sections so a worksheet can be handed out
% without the answers, then the mark scheme section added/printed
% separately.
%
% % IGCSE 0478 shared preamble -- structure only, NO topic content.
% \input this file as the FIRST line of every 05_output/*.tex (it includes
% \documentclass, so it must come before anything else). Defines the three
% custom environments every generated document must use exclusively.
% See CLAUDE.md "Strict Separation of Structure and Content".
\documentclass[11pt]{article}

\usepackage[a4paper, margin=2cm]{geometry}
\usepackage{xcolor}
\usepackage{tcolorbox}
\usepackage{enumitem}
\usepackage{amsmath}
\usepackage{hyperref}

\tcbuselibrary{skins, breakable}

% stratbox: an exam strategy / technique tip.
\newtcolorbox{stratbox}[1][Exam Strategy]{
  enhanced, breakable, colback=blue!5, colframe=blue!60!black,
  fonttitle=\bfseries, title=#1
}

% critbox: a common mistake / misconception warning.
\newtcolorbox{critbox}[1][Common Mistake]{
  enhanced, breakable, colback=red!5, colframe=red!60!black,
  fonttitle=\bfseries, title=#1
}

% scenario: an applied exam-style scenario or worked example.
\newtcolorbox{scenario}[1][Scenario]{
  enhanced, breakable, colback=green!5, colframe=green!40!black,
  fonttitle=\bfseries, title=#1
}

% \begin{document}
% \title{Topic XX: <Topic Name> -- Worksheet}
% \maketitle
% \section*{Questions}
%   \begin{enumerate}[leftmargin=*]
%     \item <question text from PastPaper_<code>.md, marks in [brackets]>
%       \begin{enumerate}[label=(\alph*)]
%         \item <sub-part text, one \item per QP label like 1(a), 1(b)>
%       \end{enumerate}
%   \end{enumerate}
%   % TODO: Insufficient source content      <- used verbatim when a QP part has no text
% \newpage
% \section*{Mark Scheme}
%   <one entry per matched question label, in the same order as above>
%   % TODO: Insufficient source content      <- used verbatim when a part has no MS match
% \end{document}

% --- Weekly tutoring document skeleton --------------------------------------
% Driven by a Study Plan in 00_syllabus/*.md (e.g. 6_month_zero_to_hero_plan.md)
% naming one week's Focus/Topics row. Content for each section is pulled from
% the Topic_XX_Name.md file(s) that row maps to (per RULE 3 mapping, confirmed
% against the syllabus). Section order is fixed; do not reorder.
%
% % IGCSE 0478 shared preamble -- structure only, NO topic content.
% \input this file as the FIRST line of every 05_output/*.tex (it includes
% \documentclass, so it must come before anything else). Defines the three
% custom environments every generated document must use exclusively.
% See CLAUDE.md "Strict Separation of Structure and Content".
\documentclass[11pt]{article}

\usepackage[a4paper, margin=2cm]{geometry}
\usepackage{xcolor}
\usepackage{tcolorbox}
\usepackage{enumitem}
\usepackage{amsmath}
\usepackage{hyperref}

\tcbuselibrary{skins, breakable}

% stratbox: an exam strategy / technique tip.
\newtcolorbox{stratbox}[1][Exam Strategy]{
  enhanced, breakable, colback=blue!5, colframe=blue!60!black,
  fonttitle=\bfseries, title=#1
}

% critbox: a common mistake / misconception warning.
\newtcolorbox{critbox}[1][Common Mistake]{
  enhanced, breakable, colback=red!5, colframe=red!60!black,
  fonttitle=\bfseries, title=#1
}

% scenario: an applied exam-style scenario or worked example.
\newtcolorbox{scenario}[1][Scenario]{
  enhanced, breakable, colback=green!5, colframe=green!40!black,
  fonttitle=\bfseries, title=#1
}

% \begin{document}
% \title{Week NN: <Focus>}
% \maketitle
% \begin{learningoutcomes}
%   <syllabus objective(s) this week's Topic_XX_Name.md content addresses>
% \end{learningoutcomes}
% \section*{Comprehensive Notes}
%   <content from 03_knowledge_base, as prose/itemize/tables; a TikZ diagram
%    only where the KB content is itself diagrammatic, e.g. hardware/network
%    layout -- never invent a diagram for content the KB describes in prose>
%   \begin{stratbox} ... \end{stratbox}    % optional, only if KB has a tip
%   \begin{critbox} ... \end{critbox}      % optional, only if KB flags a misconception
%   \begin{errortrap} ... \end{errortrap}  % optional, only if KB flags a narrow working error
% \section*{In-Class Worksheet}
%   <multi-part scenario questions built from the week's PastPaper_<code>.md
%    entries where available; mark totals come from the MS, never invented>
% \section*{Homework Assignment}
%   <one synoptic question, same sourcing rule as above>
% \newpage
% \section*{Mark Scheme}
%   \begin{modelans} ... \end{modelans}    % marker's thought process, per question
%   <longtable of marking points, sourced from the matched PastPaper_<code>.md MS>
%   % TODO: Insufficient source content      <- used verbatim when a part has no MS match
% \end{document}

\begin{document}

\title{\textbf{Topic 00: Blueprint Demo}}\\[0.4em]  
\author{IGCSE Computer Science 0478 \textbar{} SenpaiCorner \textbar{} Mr. Wei}
\date{}
\maketitle

\section{Blueprint Demo}

\subsection{Demo Subsection}

This blueprint file is not real content -- it exists only to prove the three
custom environments compile and render distinctly.

\begin{stratbox}
Use this box for an exam technique or strategy pulled from the knowledge base.
\end{stratbox}

\begin{critbox}
Use this box for a common mistake or misconception pulled from the knowledge base.
\end{critbox}

\begin{scenario}
Use this box for an applied scenario or worked example pulled from the knowledge base.
\end{scenario}

\begin{learningoutcomes}
Use this box to list the syllabus learning objective(s) a weekly document addresses.
\end{learningoutcomes}

\begin{errortrap}
Use this box for a narrow, specific working error pulled from the knowledge base (distinct from critbox's broader conceptual misconception).
\end{errortrap}

\begin{modelans}
Use this box for a marker's worked thought process alongside a mark scheme entry.
\end{modelans}

\end{document}
